\chapter{Order-Book Features}\label{ch:rs:order-book-features}

Three lots rest at the best bid and forty at the best ask, the spread is one tick. Which way does the next mid-price change
go? The book's intraday simulator answers from three hours of its own data: when the bid queue is that much smaller, the next
move is down three times in four. The order book is the most immediate information a market publishes, and every trading firm
that acts at horizons of seconds reads it. This chapter defines the standard order-book features (imbalance at the touch and
beyond, \omterm{def:rs:order-book-features:ofi}{order-flow imbalance}, the weighted mid and a \omterm{def:rs:order-book-features:micro}{microprice}, queue depletion and \omterm{def:rs:order-book-features:depletion}{cancellation rates}), measures them on
firm.tape, compares the imbalance with the prediction of a queueing model, and builds a streaming feature engine in three
languages. Its data are three hours of simulated messages: 239\,528 of them.

\section{Imbalance at the touch}

\begin{definition}[Queue imbalance]\label{def:rs:order-book-features:imbalance}
The \emph{queue imbalance}\index{queue imbalance} at the touch is $I = (q^b - q^a)/(q^b + q^a)$, where $q^b$ and $q^a$ are the
quantities resting at the best bid and the best ask; it lies in $[-1, 1]$ and is positive when the bid queue is the larger.
\end{definition}

\begin{omfigure}
\begin{tikzpicture}[font=\footnotesize,x=0.075cm,y=0.62cm]
\draw[black!60] (0,-0.5) -- (0,5.5);
\foreach \y/\p/\q in {4/100.02/35,3/100.01/28,2/100.00/40} {
  \fill[omThm!40] (0,\y-0.35) rectangle (\q,\y+0.35);
  \node[anchor=east] at (-2,\y) {\p};
  \node[anchor=west,omThm] at (\q+1,\y) {\q};
}
\foreach \y/\p/\q in {1/99.99/3,0/99.98/22,-1/99.97/31} {
  \fill[omDef!45] (0,\y-0.35) rectangle (\q,\y+0.35);
  \node[anchor=east] at (-2,\y) {\p};
  \node[anchor=west,omDef] at (\q+1,\y) {\q};
}
\node[anchor=west,omThm] at (48,3.4) {asks (lots)};
\node[anchor=west,omDef] at (48,0.2) {bids (lots)};
\draw[black!70,dashed] (-14,1.5) -- (46,1.5);
\node[anchor=west,black!70] at (48,1.5) {spread: one tick};
\node[anchor=west,align=left] at (48,-1.0) {touch imbalance $I = (3 - 40)/43 = -0.86$};
\end{tikzpicture}

\omcaption{The book around the touch: three lots at the best bid (99.99), forty at the best ask (100.00). The \omterm{def:rs:order-book-features:imbalance}{queue imbalance} at the
touch is $-0.86$; over three levels on each side it is $(56 - 103)/159 = -0.30$, because the deeper bids are larger. A schematic, not
data.}\label{fig:rs:order-book-features:book}
\end{omfigure}

When the spread is one tick, the mid moves only when one of the two best queues empties (or a new level appears inside a wider
spread). A small bid queue is closer to emptying than a large ask queue, so imbalance predicts the direction of the next move.
The first model of this is a race.

\begin{proposition}[Imbalance as a race of two queues]\label{prop:rs:order-book-features:race}
Model the two best queues, in lots, as independent birth--death processes with arrival rate $\lambda$ and departure rate $\mu$
(Book 4, chapter 8), starting from $q^b$ and $q^a$. The probability that the next mid move is up is the probability that the ask
queue empties first, $\int_0^\infty(1 - F_{q^b}(t))\,dF_{q^a}(t)$, where $F_q$ is the law of the time for a queue starting at $q$ to
empty. It increases with $q^b$, decreases with $q^a$, and equals $\tfrac12$ when $q^b = q^a$.
\end{proposition}

\begin{proof}
The next move is up when the ask side empties before the bid side; with independent queues this is the stated race. A larger
starting queue takes stochastically longer to empty, which gives the monotonicity; symmetry gives $\tfrac12$.
\end{proof}

On the simulator's three hours, 2.03 lots a second arrive at each best queue and 2.13 leave it (additions at the best price;
cancellations and executions there). \Cref{fig:rs:order-book-features:updown} sets the race's prediction against the data, bucket
by bucket of the imbalance. The direction is right: the probability of an up move rises from 0.245 in the most ask-heavy tenth of
observations to 0.736 in the most bid-heavy. The model is far too confident: it predicts 0.017 and 0.979. The reason is the
mechanism the simulator plants, and that real markets share: the queues are not independent random walks, because the
efficient price moves and the liquidity providers on its wrong side cancel. Imbalance is a noisy view of that, and a queueing
model that ignores it overstates what the view is worth. Gould and Bonart (2016) fit the same relation by logistic regression on
ten Nasdaq stocks and find a strongly significant improvement over a null model, large for large-tick stocks.

\begin{omfigure}
\begin{tikzpicture}
\begin{axis}[width=10cm,height=5.4cm,xlabel={queue imbalance $I$ (bucket centre)},ylabel={P(next mid move is up)},
  tick label style={font=\small},label style={font=\small},xmin=-1,xmax=1,ymin=0,ymax=1,grid=major,grid style={gray!25},
  legend columns=2,legend style={font=\footnotesize,at={(0.5,-0.32)},anchor=north,draw=none,fill=none,
  /tikz/every even column/.append style={column sep=8pt}},table/col sep=comma]
\addplot[omDef,very thick,mark=*] table[x=centre,y=empirical]{figdata/research/08-order-book-features/updown.csv};
\addplot[omThm,thick,dashed,mark=square] table[x=centre,y=model]{figdata/research/08-order-book-features/updown.csv};
\addplot[black!50,thin,sharp plot] coordinates {(-1,0.5) (1,0.5)};
\legend{simulated tape,race of two birth--death queues}
\end{axis}
\end{tikzpicture}

\omcaption{Probability that the next mid-price move is up, by tenth of \omterm{def:rs:order-book-features:imbalance}{queue imbalance} at a one-tick spread: 0.245 to 0.736 in three
hours of \texttt{firm.tape}, against 0.017 to 0.979 for the race of two independent birth--death queues with the rates measured on the
same data. Data: \texttt{firm.tape}, seed 8; \texttt{firm.queues}.}\label{fig:rs:order-book-features:updown}
\end{omfigure}

\section{Depth beyond the touch}

\begin{definition}[Depth imbalance, depth profile]\label{def:rs:order-book-features:depth}
The \emph{depth imbalance}\index{depth imbalance} over $L$ levels is $(\sum_{i \le L}q^b_i - \sum_{i \le L}q^a_i)/(\sum_{i\le L}q^b_i
+ \sum_{i \le L}q^a_i)$, with $q^b_i, q^a_i$ the quantities at the $i$-th best bid and ask prices. The \emph{depth profile}\index{depth
profile} of a side is the vector of its quantities by distance from the touch.
\end{definition}

Depth beyond the touch is where the next best quote will come from once a queue empties, and where large orders show their hand, or
hide it. Its information is slower than the touch's and more easily faked (orders far from the touch are cheap to place and cancel,
which Book~9, chapter~29, returns to). A \omterm{def:rs:order-book-features:depth}{depth profile} is a feature vector rather than a number: its slope, its concentration, the
distance to the first level holding more than a given quantity. The engine of this chapter computes the imbalance over the first five
levels; the \omterm{def:rs:order-book-features:depth}{depth profile} itself is the input of the machine-learned models of Book~12.

\section{Order-flow imbalance}

\begin{definition}[Order-flow imbalance]\label{def:rs:order-book-features:ofi}
Let $b_n, q^b_n, a_n, q^a_n$ be the best bid, its size, the best ask and its size after the $n$-th book event. The \emph{order-flow
imbalance}\index{order-flow imbalance} contribution of event $n$ is
\[
e_n = \mathbf 1_{b_n \ge b_{n-1}}q^b_n - \mathbf 1_{b_n \le b_{n-1}}q^b_{n-1} - \mathbf 1_{a_n \le a_{n-1}}q^a_n + \mathbf 1_{a_n \ge
a_{n-1}}q^a_{n-1},
\]
and the order-flow imbalance (OFI) over an interval is the sum of the $e_n$ in it (Cont, Kukanov and Stoikov, 2014).
\end{definition}

The four terms count demand added at the bid, demand removed from it, supply added at the ask and supply removed from it; when a
price level changes, the whole old queue counts as removed and the whole new one as added. Cont, Kukanov and Stoikov found on 50 US
stocks that over short intervals price changes are mainly driven by OFI, through a linear relation whose slope is inversely proportional
to market depth.

\begin{method}[Measuring the impact of order flow]\label{met:rs:order-book-features:ofi}
On a grid of intervals of length $\Delta$, regress the change in mid price (in ticks) on the interval's OFI divided by the average size
of the best queues; the slope is the price impact per unit of depth-normalised flow, and the $R^2$ says how much of the price change the
flow explains at that horizon.
\end{method}

On the simulated tape the regression explains 10\% of mid changes at one second, 53\% at ten seconds and 79\% at sixty; the slope at ten
seconds is 0.47 ticks per average best-queue size of net flow (\cref{fig:rs:order-book-features:ofi}). At one second most intervals have
no mid change at all, and those that do are single ticks triggered by the last event, which OFI sees only in part.

\begin{omfigure}
\begin{tikzpicture}
\begin{axis}[width=9.5cm,height=5.2cm,xlabel={OFI over 10 seconds (in average best-queue sizes)},ylabel={mid change (ticks)},
  tick label style={font=\small},label style={font=\small},grid=major,grid style={gray!25},table/col sep=comma]
\addplot[omDef,only marks,mark=*,mark size=1.8pt] table[x=ofi,y=dmid]{figdata/research/08-order-book-features/ofi.csv};
\addplot[omThm,thick,domain=-3.5:3.5,samples=2] {0.4664*x};
\end{axis}
\end{tikzpicture}

\omcaption{Mid-price change over ten-second intervals against the interval's \omterm{def:rs:order-book-features:ofi}{order-flow imbalance}, in twenty bins of equal count, with
the regression line (slope 0.47 ticks per unit, $R^2 = 0.53$ on the 1\,078 intervals). Data: \texttt{firm.tape} through
\texttt{firm.lobfeat}, seed 8.}\label{fig:rs:order-book-features:ofi}
\end{omfigure}

\section{The microprice}

\begin{definition}[Weighted mid price, microprice]\label{def:rs:order-book-features:micro}
The \emph{weighted mid price}\index{weighted mid price} is $(a\,q^b + b\,q^a)/(q^b + q^a)$: the mid moved towards the side with the
smaller queue. The book's \emph{microprice}\index{microprice} is the mid plus the expected size and sign of the next mid change given the
current imbalance bucket (at a one-tick spread), estimated on past data.
\end{definition}

The weighted mid uses the imbalance through a fixed formula; the \omterm{def:rs:order-book-features:micro}{microprice} learns the adjustment from data. Stoikov (2018) builds the full
estimator from the whole sequence of later mid changes and finds it a better \omterm{def:rs:anatomy-of-a-predictor:predictor}{predictor} of short-term prices than either the mid or the
weighted mid; the book's version keeps only the first term of that sequence. Fitted on the first half of the simulated session and scored
on the second, it has a mean squared error 10\% below the mid's for the mid one second later; the weighted mid is 27\% worse at one second
but the best at five and thirty seconds (12\% and 3\% below the mid). A first-order correction is right about the next move and blind
beyond it; the weighted mid is crude but keeps pointing the right way.

\begin{omfigure}
\begin{tikzpicture}
\begin{axis}[width=8.5cm,height=5cm,ybar,bar width=14pt,ylabel={MSE relative to the mid},tick label style={font=\small},
  label style={font=\small},xmin=-0.6,xmax=2.6,ymin=0,ymax=1.4,xtick={0,1,2},xticklabels={1 second,5 seconds,30 seconds},
  grid=major,grid style={gray!25},legend columns=2,legend style={font=\footnotesize,at={(0.5,-0.25)},anchor=north,draw=none,
  fill=none,/tikz/every even column/.append style={column sep=8pt}},table/col sep=comma]
\addplot[fill=omThm!45,draw=omThm] table[x=k,y=wmid]{figdata/research/08-order-book-features/forecast.csv};
\addplot[fill=omDef!55,draw=omDef] table[x=k,y=micro]{figdata/research/08-order-book-features/forecast.csv};
\addplot[black!70,thin,sharp plot,forget plot] coordinates {(-0.6,1) (2.6,1)};
\legend{weighted mid,microprice (one step)}
\end{axis}
\end{tikzpicture}

\omcaption{Mean squared error of the weighted mid and of the one-step \omterm{def:rs:order-book-features:micro}{microprice} as forecasts of the mid 1, 5 and 30 seconds later,
relative to the mid's own error, on the second half of the simulated session (the \omterm{def:rs:order-book-features:micro}{microprice}'s table is fitted on the first). Data:
\texttt{firm.tape} through \texttt{firm.lobfeat}, seed 8.}\label{fig:rs:order-book-features:forecast}
\end{omfigure}

\section{Queue depletion and cancellations}

\begin{definition}[Queue depletion rate, cancellation rate]\label{def:rs:order-book-features:depletion}
The \emph{queue depletion rate}\index{queue depletion rate} of a best queue is the quantity leaving it per second (cancellations and
executions) over a trailing window; its \emph{cancellation rate}\index{cancellation rate} counts cancellations only.
\end{definition}

Rates carry what sizes do not: a queue of forty lots losing ten a second is weaker than one of twenty losing none. In the simulator,
liquidity providers cancel faster on the side the efficient price has moved away from, as real ones do when their quotes become stale.
The difference between the cancellations at the best ask and at the best bid over the last second is therefore a direct view of the
hidden price: its correlation with the direction of the next mid move is 0.30, against 0.25 for the \omterm{def:rs:order-book-features:imbalance}{queue imbalance}. The general lesson
is the chapter's: the most informative features are those closest to the mechanism that moves the price.

\section{Predictor cards}

\begin{predictorcard}{Queue imbalance}\label{pred:rs:order-book-features:imbalance}
\sfield{Definition} $I = (q^b - q^a)/(q^b + q^a)$ at the touch, after each message.
\sfield{Inputs and timestamps} The book after the last message received; the receive time is the \omterm{def:rs:point-in-time-data-and-the-biases:bitemporal}{knowledge time}.
\sfield{Rationale} The smaller queue empties first (\cref{prop:rs:order-book-features:race}); liquidity providers on the wrong side of the
efficient price withdraw.
\sfield{Horizon and half-life} The next mid move, seconds. P(up) from 0.245 to 0.736 across tenths of $I$ on the simulated tape.
\sfield{Normalisation} None needed: bounded; bucketed or fed to a logistic model.
\sfield{Failure modes} Spoofed size (Book~9, chapter~29); hidden and iceberg orders; wide spreads, where moves need not empty a queue.
\sfield{Sources} Gould and Bonart (2016); \texttt{rs\_orderbook.up\_probability}.
\end{predictorcard}

\begin{predictorcard}{Order-flow imbalance}\label{pred:rs:order-book-features:ofi}
\sfield{Definition} The sum of the $e_n$ over the last $\Delta$ seconds, divided by the average best-queue size.
\sfield{Inputs and timestamps} Every book event, in sequence, with receive times.
\sfield{Rationale} Net demand at the touch moves the price linearly, with a slope inversely proportional to depth.
\sfield{Horizon and half-life} Contemporaneous in its definition; as a \omterm{def:rs:anatomy-of-a-predictor:predictor}{predictor}, the next seconds. $R^2$ 10\%, 53\% and 79\% at 1, 10 and 60
seconds (contemporaneous) on the simulated tape.
\sfield{Normalisation} By depth; by volatility across instruments.
\sfield{Failure modes} Needs every event in order (a gap in the feed corrupts the sum); spoofed additions count as demand.
\sfield{Sources} Cont, Kukanov and Stoikov (2014); \texttt{rs\_orderbook.ofi\_regression}.
\end{predictorcard}

\begin{predictorcard}{Stale-quote cancellation imbalance}\label{pred:rs:order-book-features:cancel}
\sfield{Definition} Lots cancelled at the best ask minus lots cancelled at the best bid over the last second.
\sfield{Inputs and timestamps} Cancel messages with their prices and receive times.
\sfield{Rationale} Liquidity providers pull quotes the efficient price has left behind.
\sfield{Horizon and half-life} The next mid move; correlation 0.30 with its direction on the simulated tape.
\sfield{Normalisation} By the queues' sizes, or ranked across instruments.
\sfield{Failure modes} Cancel-and-replace at the same price by an algorithm that refreshes its quotes; venue differences in how amendments
are reported.
\sfield{Sources} The simulator's mechanism (chapter~2); \texttt{rs\_orderbook.cancel\_feature}.
\end{predictorcard}

\section{Tutorial: reading the book}\label{tut:rs:order-book-features:read}

\begin{tutorial}
\textbf{Goal.} Run three hours of simulated messages through the feature engine, measure imbalance, OFI and the \omterm{def:rs:order-book-features:micro}{microprice}, and compare the
imbalance with the queueing model. \textbf{End state:}
\cref{fig:rs:order-book-features:updown,fig:rs:order-book-features:ofi,fig:rs:order-book-features:forecast}.
\begin{tutsteps}
\item \textbf{The engine.} Apply the message to the order and level maps, read the new top, compute the OFI contribution from the previous
top, then the imbalances, the weighted mid and the \omterm{def:rs:order-book-features:micro}{microprice}.
\omcode{code/firm/lobfeat/firm_lobfeat.py}{66}{96}{One message through the Python reference engine.}\label{lst:rs:order-book-features:engine}
\item \textbf{The race.} The probability that the ask queue empties first, from the depletion laws of Book 4's \texttt{firm.queues},
conditioned on one of them emptying within ten minutes.
\omcode{code/research/08-order-book-features/python/rs_orderbook.py}{79}{86}{The birth--death race.}\label{lst:rs:order-book-features:race}
\item \textbf{Run} \texttt{up\_probability()}, \texttt{best\_level\_rates()}, \texttt{ofi\_regression} at 1, 10 and 60 seconds,
\texttt{forecast\_errors()}, \texttt{cancel\_feature()} and \texttt{fig\_orderbook.py}.
\end{tutsteps}
\textbf{What to change next.} Fit a logistic regression of the next move's direction on $I$ and on the cancellation imbalance together;
estimate the \omterm{def:rs:order-book-features:micro}{microprice} table on spreads of two ticks as well.
\end{tutorial}

\section{Build: the streaming feature engine}\label{bld:rs:order-book-features:lobfeat}

\begin{build}
\bfield{Purpose} The order-book features of the miniature firm, computed message by message in the same way in research (Python) and in
production (C++20, with a Rust twin); Book~11's market maker and Book~12's order-book model consume them.
\bfield{Interface} \texttt{Engine(levels, g).on(kind, oid, side, price, qty)} returning \texttt{Features} (best quotes and sizes,
\texttt{imbalance}, \texttt{depth\_imbalance}, \texttt{ofi}, \texttt{ofi\_cum}, \texttt{wmid}, \texttt{micro}); \texttt{run(msgs)};
\texttt{LobFeatures} in \texttt{cpp/firm\_lobfeat.hpp} and \texttt{rust/src/lib.rs}.
\bfield{Rules} Integer prices and sizes; no features until both sides have a quote; the OFI of the first two-sided event is zero; identical
arithmetic in the three languages.
\bfield{Acceptance tests} \texttt{code/firm/lobfeat/tests/}: a hand-checked sequence of adds, executions and a level change; the committed
fixture reproduced by the Python reference; bounded imbalances and OFI tracking the mid; the C++ and Rust tests replay the 2\,769-message
fixture and match every feature of the 2\,766 two-sided rows.
\bfield{Stretch} \omterm{def:rs:order-book-features:depth}{Depth profiles} as fixed-size vectors; trailing-window depletion and \omterm{def:rs:order-book-features:depletion}{cancellation rates} in the engine; a lock-free
single-writer queue feeding it (Book~13).
\end{build}

The C++ core, the one subtle part of the engine, is the order-flow update from the previous top of book:

\omcode{code/firm/lobfeat/cpp/firm_lobfeat.hpp}{43}{60}{The C++20 engine after each message.}\label{lst:rs:order-book-features:cpp}

\begin{omsources}
\item R.~Cont, A.~Kukanov and S.~Stoikov, ``The price impact of order book events'', \emph{Journal of Financial Econometrics} 12(1), 2014.
\item S.~Stoikov, ``The micro-price: a high-frequency estimator of future prices'', \emph{Quantitative Finance} 18(12), 2018.
\item M.~D.~Gould and J.~Bonart, ``Queue imbalance as a one-tick-ahead price predictor in a limit order book'', \emph{Market Microstructure
and Liquidity} 2(2), 2016.
\item R.~Cont, S.~Stoikov and R.~Talreja, ``A stochastic model for order book dynamics'', \emph{Operations Research} 58(3), 2010.
\item W.~Huang, C.-A.~Lehalle and M.~Rosenbaum, ``Simulating and analyzing order book data: the queue-reactive model'', \emph{Journal of the
American Statistical Association} 110(509), 2015.
\end{omsources}

\section{Exercises}

\begin{exercise}[$\star$]\label{exo:rs:order-book-features:1}
Best bid 99.98 for 700 shares, best ask 99.99 for 300. What are the \omterm{def:rs:order-book-features:imbalance}{queue imbalance}, the mid and the weighted mid?
\end{exercise}

\begin{exercise}[$\star$]\label{exo:rs:order-book-features:2}
The bid moves from 99.98 (500 shares) to 99.99 (200 shares) and the ask stays at 100.00 (400 shares, unchanged). What is the event's OFI
contribution?
\end{exercise}

\begin{exercise}[$\star$]\label{exo:rs:order-book-features:3}
The OFI slope is 0.47 ticks per average best-queue size, and that size is 2\,750 shares. What mid move does a net buying flow of 11\,000
shares over ten seconds predict?
\end{exercise}

\begin{exercise}[$\star\star$]\label{exo:rs:order-book-features:4}
For the race with $\lambda = \mu$ and queues of one lot at the bid and $n$ lots at the ask, argue that the probability of an up move tends to
zero as $n$ grows. What does that say about the model's extreme buckets?
\end{exercise}

\begin{exercise}[$\star\star$]\label{exo:rs:order-book-features:5}
Why does OFI's $R^2$ rise with the length of the interval, from 10\% at one second to 79\% at sixty?
\end{exercise}

\begin{exercise}[$\star\star$]\label{exo:rs:order-book-features:6}
Show that the weighted mid equals the mid plus $\tfrac s2 I$, with $s$ the spread and $I$ the \omterm{def:rs:order-book-features:imbalance}{queue imbalance}.
\end{exercise}

\begin{exercise}[$\star\star\star$]\label{exo:rs:order-book-features:7}
\emph{Coding.} Using the tape of the chapter, compute the correlation of the \omterm{def:rs:order-book-features:depth}{depth imbalance} over five levels with the direction of the
next mid move, and compare it with the touch imbalance's.
\end{exercise}

\begin{exercise}[$\star\star\star$]\label{exo:rs:order-book-features:8}
\emph{Find the flaw.} ``Our order-book model predicts the next mid move with 74\% accuracy. We computed the features from the consolidated
book stamped with exchange timestamps and the target from the same feed.''
\end{exercise}

\section{Problem: Three Lots Against Forty}

\begin{problem}[{Weekend problem --- how much does the order book know, and how much does a queueing model think it knows?}]\label{pb:rs:order-book-features:1}
Three hours of \texttt{firm.tape} (seed 8), 239\,528 messages, through \texttt{firm.lobfeat}.

\medskip\noindent\textbf{Part I --- Imbalance.}
\begin{enumerate}
\item What are the probabilities of an up move in the lowest and highest tenths of the \omterm{def:rs:order-book-features:imbalance}{queue imbalance}?
\item How many observations fall in each of those two tenths?
\item What are the average queue sizes (lots) in the extreme tenths?
\item At what imbalance is the probability one half?
\end{enumerate}

\medskip\noindent\textbf{Part II --- The model.}
\begin{enumerate}[resume]
\item What arrival and departure rates (lots per second) does each best queue have?
\item What does the race predict for the extreme tenths?
\item Why is the model too confident?
\item Which mechanism of the simulator would a correct model have to include?
\end{enumerate}

\medskip\noindent\textbf{Part III --- Flow.}
\begin{enumerate}[resume]
\item What are OFI's $R^2$ at 1, 10 and 60 seconds, and its slope at ten seconds?
\item What is the average best-queue size used to normalise it?
\item How does the one-step \omterm{def:rs:order-book-features:micro}{microprice} compare with the mid at one second?
\item Which forecast wins at five and thirty seconds, and by how much?
\item What is the cancellation imbalance's correlation with the next move, and the \omterm{def:rs:order-book-features:imbalance}{queue imbalance}'s?
\end{enumerate}

\medskip\noindent\textbf{Part IV --- Engineering and judgement.}
\begin{enumerate}[resume]
\item Why must the three implementations of the engine produce identical numbers?
\item How many fixture rows do the C++ and Rust tests check?
\item Which features need every message in sequence, and which survive a snapshot feed?
\item Which of the chapter's features would a spoofer move most easily?
\item State the \emph{named result}: the empirical and modelled probabilities of an up move in the extreme tenths of imbalance, and OFI's
$R^2$ at ten seconds.
\item What would you add to the race model first?
\item In one sentence: what does the order book tell you?
\end{enumerate}
\end{problem}

\section{Interview questions}

\begin{interviewq}[$\star$ \iqroles{trader, researcher}]\label{iq:rs:order-book-features:1}
The bid queue is ten times the ask queue at a one-tick spread. What do you expect the next price move to be, and why?
\end{interviewq}

\begin{interviewq}[$\star\star$ \iqroles{researcher}]\label{iq:rs:order-book-features:2}
Define \omterm{def:rs:order-book-features:ofi}{order-flow imbalance} and explain why it relates linearly to price changes.
\end{interviewq}

\begin{interviewq}[$\star\star$ \iqroles{researcher, developer}]\label{iq:rs:order-book-features:3}
What is a \omterm{def:rs:order-book-features:micro}{microprice}? How would you estimate one?
\end{interviewq}

\begin{interviewq}[$\star\star$ \iqroles{developer}]\label{iq:rs:order-book-features:4}
Your research features are in Python and production is in C++. How do you guarantee they compute the same thing?
\end{interviewq}

\begin{interviewq}[$\star\star$ \iqroles{trader, researcher}]\label{iq:rs:order-book-features:5}
How can an order-book feature be manipulated, and how would you make a model robust to it?
\end{interviewq}

\begin{interviewq}[$\star\star\star$ \iqroles{researcher}]\label{iq:rs:order-book-features:6}
Model the best bid and ask queues as independent birth--death processes. What does the model predict for the next move, and what does it
leave out?
\end{interviewq}
